\documentclass[11pt,a4paper]{article}

% ---------- Encoding e idioma ----------
\usepackage[utf8]{inputenc}
\usepackage[T1]{fontenc}
\usepackage[spanish,es-tabla]{babel}

% ---------- Diseño de página ----------
\usepackage[margin=2.5cm]{geometry}
\usepackage{setspace}
\onehalfspacing

% ---------- Matemáticas ----------
\usepackage{amsmath,amssymb,amsthm,mathtools}

% ---------- Figuras y tablas ----------
\usepackage{graphicx}
\usepackage{booktabs}
\usepackage{caption}
\usepackage{subcaption}
\usepackage{float}

% ---------- Código fuente ----------
\usepackage{xcolor}
\usepackage{listings}

\definecolor{codegray}{rgb}{0.95,0.95,0.95}
\definecolor{codekey}{rgb}{0.10,0.30,0.70}
\definecolor{codecom}{rgb}{0.20,0.55,0.20}
\definecolor{codestr}{rgb}{0.60,0.10,0.10}

\lstdefinestyle{matlab}{
    language=Matlab,
    backgroundcolor=\color{codegray},
    basicstyle=\ttfamily\footnotesize,
    keywordstyle=\color{codekey}\bfseries,
    commentstyle=\color{codecom}\itshape,
    stringstyle=\color{codestr},
    numbers=left,
    numberstyle=\tiny\color{gray},
    numbersep=8pt,
    frame=single,
    rulecolor=\color{black!30},
    breaklines=true,
    breakatwhitespace=true,
    showstringspaces=false,
    tabsize=2,
    captionpos=b,
    escapeinside={(*@}{@*)}
}
\lstset{style=matlab}

% ---------- Enlaces ----------
\usepackage[colorlinks=true, linkcolor=blue!60!black, urlcolor=blue!60!black, citecolor=green!50!black]{hyperref}

% ---------- Metadatos ----------
\title{
    \vspace{-1.5cm}
    \textbf{Proyecto de Visión por Computadora (1MTR27)} \\[0.3cm]
    \Large Parte 1: Algoritmo Automático de Preprocesamiento \\ de Imágenes de Escritorio
}
\author{Jamin Yauri Cajas \\ \texttt{jaminyauricajas@gmail.com}}
\date{24 de abril de 2026}

% ============================================================
\begin{document}
\maketitle

\begin{abstract}
\noindent
Este informe presenta un algoritmo completamente automático para el preprocesamiento
de imágenes de escritorio con objetos de uso diario (lápices, bolígrafos, marcadores,
cuadernos, calculadoras). El pipeline combina filtrado no lineal para remover ruido
impulsivo, filtrado gaussiano suave para ruido residual, balance de blancos
\emph{gray-world} para corregir dominantes cromáticas de la iluminación, y ecualización
adaptativa del histograma con restricción de contraste (CLAHE) aplicada sobre el canal
de luminancia del espacio CIE Lab para normalizar la intensidad sin alterar la
cromaticidad original. Todos los pasos usan parámetros universales que no requieren
ajuste manual por imagen; la adaptación ocurre dentro de cada operador. Se evalúa el
pipeline sobre un conjunto de 17 imágenes tomadas bajo condiciones de iluminación
heterogéneas, reportando PSNR y SSIM contra la imagen original como medida del grado
de modificación aplicado.
\end{abstract}

% ============================================================
\section{Introducción}
\label{sec:intro}

El enunciado del proyecto exige, como paso previo a la detección y reconocimiento
de objetos sobre un escritorio, un algoritmo de preprocesamiento que:
\begin{enumerate}
    \item Elimine ruidos y manchas presentes en la imagen capturada.
    \item Modifique las intensidades para favorecer la detección posterior.
    \item Incluya cualquier paso adicional que se considere necesario, debidamente justificado.
    \item Sea \textbf{completamente automático}: la entrada es una imagen y el proceso
          no debe requerir ajuste de parámetros caso por caso.
\end{enumerate}

La dificultad principal de este requisito de automatización es que las imágenes se
toman bajo condiciones no controladas: iluminación de intensidad y temperatura de color
variables, ángulos y distancias arbitrarias, ruido dependiente del sensor de cada
dispositivo. Cualquier paso del pipeline debe, por tanto, adaptarse al contenido
observado en lugar de depender de un ajuste global.

La solución propuesta se implementa en la función
\lstinline{pipeline_preproc.m} (ver
\lstinline{src/preprocessing/pipeline_preproc.m}) y se organiza en cinco etapas
secuenciales descritas en la Sección~\ref{sec:pipeline}.

% ============================================================
\section{Metodología}
\label{sec:pipeline}

\subsection{Vista general del pipeline}

El pipeline sigue el orden que se resume en la Tabla~\ref{tab:pipeline}.

\begin{table}[H]
\centering
\caption{Etapas del pipeline de preprocesamiento.}
\label{tab:pipeline}
\begin{tabular}{@{}clll@{}}
\toprule
\# & Etapa & Operador & Requisito del enunciado \\
\midrule
1 & Normalización de rango & \(\texttt{im2double}\) & Preparación \\
2 & Denoise impulsivo & Mediana \(3\times 3\) & Ruido/manchas \\
3 & Denoise residual & Gaussiano (\(\sigma=0{,}8\)) & Ruido/manchas \\
4 & Balance de blancos & Gray-world & Procesamiento adicional \\
5 & Normalización de intensidad & CLAHE sobre \(L\) en Lab & Modificación de intensidades \\
\bottomrule
\end{tabular}
\end{table}

El orden no es arbitrario: se justifica explícitamente en la Sección~\ref{sec:orden}.

% ----------------------------------------------------------
\subsection{Etapa 1 — Normalización a \texttt{double} en \([0,1]\)}
\label{sec:norm}

Una imagen RGB capturada con una cámara estándar se representa habitualmente como
una matriz de enteros \texttt{uint8} cuyo rango es \(\{0, 1, \dots, 255\}\). Los
operadores aritméticos sobre este tipo sufren de \textbf{saturación y truncamiento}:
por ejemplo, \(200 + 100\) se satura a \(255\) en lugar de producir \(300\), y una
división entera como \(5/2\) se trunca a \(2\). Todos los operadores posteriores del
pipeline — promedios espaciales, divisiones del balance de blancos, conversión
\(\text{RGB} \to \text{Lab}\), ecualización del histograma — requieren aritmética de
punto flotante para preservar información.

La conversión se realiza según:
\begin{equation}
    I'(x,y,c) = \frac{I(x,y,c)}{255}, \qquad I' \in [0,1]^{H\times W\times 3}.
    \label{eq:norm}
\end{equation}

\subsection{Etapa 2 — Denoise no lineal con filtro de mediana}
\label{sec:median}

Los artefactos de tipo \emph{sal y pimienta} — píxeles individuales con valor extremo,
frecuentes en fotografías con compresión JPEG agresiva o sensores CMOS económicos —
no se eliminan bien con filtros lineales. El filtro de mediana, definido por
\begin{equation}
    I_{\text{med}}(x,y,c) = \operatorname{mediana} \bigl\{ I(u,v,c) \;:\; (u,v) \in \mathcal{N}_{3\times 3}(x,y) \bigr\},
    \label{eq:median}
\end{equation}
es un operador \textbf{no lineal} y por tanto \textbf{robusto a outliers}: un píxel
anómalo en una vecindad de nueve no altera la mediana, ya que esta corresponde al
quinto valor del conjunto ordenado. En contraste con el filtro promedio, preserva
mejor los bordes porque no introduce valores intermedios inexistentes en la imagen.

La operación se aplica independientemente a cada canal de color, pues el ruido
impulsivo puede presentarse en un único canal (efecto del demosaico del Bayer pattern).

\subsection{Etapa 3 — Suavizado gaussiano}
\label{sec:gauss}

Tras la mediana persiste un \textbf{ruido gaussiano residual} de menor amplitud,
originado por la amplificación del sensor. Se suaviza con convolución por un kernel
gaussiano bidimensional:
\begin{equation}
    G_\sigma(x,y) = \frac{1}{2\pi\sigma^2}\exp\!\left(-\frac{x^2+y^2}{2\sigma^2}\right),
    \qquad I_{\text{gauss}} = I_{\text{med}} * G_\sigma.
    \label{eq:gauss}
\end{equation}
Con \(\sigma = 0{,}8\) el kernel efectivo es de \(\sim 5\times 5\) píxeles, lo que
reduce el ruido de alta frecuencia sin destruir bordes de objetos finos como lápices
o la tipografía impresa en marcadores.

La \textbf{secuencia mediana~\(\to\)~gaussiano} es deliberada: si se invirtiera, el
filtro gaussiano distribuiría la energía del outlier sobre sus vecinos, haciendo
imposible su eliminación posterior.

\subsection{Etapa 4 — Balance de blancos por la hipótesis del mundo gris}
\label{sec:wb}

La iluminación de la escena introduce un \textbf{dominante cromático}: una bombilla
incandescente tiñe toda la imagen de amarillo, un fluorescente la vuelve verdosa.
El algoritmo \emph{gray-world} asume que la media global del color en una escena
natural es aproximadamente gris neutro, y escala cada canal para imponer esa
propiedad. Sea
\begin{equation}
    \bar{c} = \frac{1}{HW}\sum_{x,y} I_{\text{gauss}}(x,y,c), \qquad c \in \{R,G,B\},
\end{equation}
y \(\bar{m} = \tfrac{1}{3}(\bar{R}+\bar{G}+\bar{B})\) la luminancia promedio.
La imagen balanceada es
\begin{equation}
    I_{\text{wb}}(x,y,c) = I_{\text{gauss}}(x,y,c) \cdot \frac{\bar{m}}{\bar{c}},
    \label{eq:grayworld}
\end{equation}
seguida de un clipping a \([0,1]\).

La principal virtud de esta formulación para el requisito de automaticidad es que
\textbf{no posee parámetros}: se auto-calibra a partir de los estadísticos de la
imagen. Su principal limitación ocurre cuando una clase cromática domina la escena
(p.\,ej., un escritorio con exclusivamente marcadores rojos); en el dataset empleado
esta situación no se presenta.

\subsection{Etapa 5 — Normalización de intensidad con CLAHE en el canal \(L\)}
\label{sec:clahe}

La \textbf{ecualización de histograma clásica} redistribuye las intensidades para
aproximar una distribución uniforme, lo que incrementa el contraste global pero:
(i) amplifica el ruido, y (ii) falla en imágenes con gradientes de iluminación, donde
una zona oscura y una clara no pueden corregirse simultáneamente con una única
transformación.

CLAHE (\emph{Contrast Limited Adaptive Histogram Equalization}) resuelve ambos
problemas. El procedimiento es:
\begin{enumerate}
    \item Particionar la imagen en una grilla de \(8\times 8\) \emph{tiles}.
    \item Para cada tile, calcular su histograma y su función de distribución acumulada (CDF).
    \item Aplicar un \textbf{clip limit} \(\beta\) al histograma: los bins con frecuencia mayor a \(\beta\) se recortan, y el excedente se redistribuye uniformemente. Esto evita amplificar ruido en regiones de baja varianza.
    \item Ecualizar cada tile usando su CDF modificada.
    \item Interpolar bilinealmente las transformaciones entre tiles vecinos para evitar artefactos de bloques.
\end{enumerate}

Formalmente, para un tile \(T\) con \(N\) píxeles y clip limit relativo \(\beta\):
\begin{equation}
    h_{\text{clip}}(k) = \min\!\bigl(h_T(k),\, \beta N\bigr),\qquad
    T_{\text{eq}}(I) = \operatorname{CDF}(h_{\text{clip}})(I).
\end{equation}

\subsubsection*{¿Por qué operar en Lab y no en RGB?}

Aplicar CLAHE directamente en cada canal RGB produce \textbf{aberraciones cromáticas}:
al alterar de manera no sincronizada la intensidad de \(R\), \(G\) y \(B\), se desplaza
el color percibido de cada píxel. El espacio \textbf{CIE Lab} separa la luminancia
\(L\) de la cromaticidad \((a, b)\), permitiendo modificar el contraste sin tocar el
tono. La transformación es
\begin{equation}
    (R,G,B) \xrightarrow{\text{rgb2lab}} (L,a,b),\qquad
    L' = \operatorname{CLAHE}(L),\qquad
    (L',a,b) \xrightarrow{\text{lab2rgb}} (R',G',B').
    \label{eq:lab}
\end{equation}

\subsection{Clamping final y salida}

Tras la transformación inversa \(\text{Lab}\to\text{RGB}\), los valores pueden caer
ligeramente fuera de \([0,1]\) por redondeo numérico. Se aplica un clipping final:
\begin{equation}
    I_{\text{out}}(x,y,c) = \max\bigl(0,\min(1, I(x,y,c))\bigr).
\end{equation}

% ============================================================
\section{Justificación del orden del pipeline}
\label{sec:orden}

El orden \texttt{normalización} \(\to\) \texttt{mediana} \(\to\) \texttt{gaussiano}
\(\to\) \texttt{gray-world} \(\to\) \texttt{CLAHE} responde a dependencias causales
entre etapas:

\begin{enumerate}
    \item La \textbf{normalización} precede a todo: habilita la aritmética precisa.
    \item La \textbf{mediana} precede al gaussiano porque este último ``unta'' los
          outliers sobre sus vecinos, haciendo imposible eliminarlos después.
    \item El \textbf{balance de blancos} precede a CLAHE porque un tinte cromático
          global perturbaría el histograma del canal \(L\); eliminarlo primero
          permite que CLAHE trabaje sobre una luminancia limpia.
    \item El \textbf{CLAHE sobre \(L\)} se ejecuta último porque es el paso de
          realce de contraste, que debe aplicarse sobre una imagen ya libre de ruido
          y con color balanceado.
\end{enumerate}

% ============================================================
\section{Carácter automático del algoritmo}
\label{sec:auto}

El enunciado exige que el pipeline funcione sin ajuste manual de parámetros por
imagen. Esto se logra mediante una observación crucial: los parámetros del código
\emph{no necesitan} adaptarse por imagen porque los propios operadores se
auto-adaptan al contenido. En particular:

\begin{itemize}
    \item CLAHE se adapta al histograma local de cada tile de cada imagen.
    \item \emph{Gray-world} se adapta a los estadísticos cromáticos de cada imagen.
    \item El filtro de mediana es, por construcción, insensible a la amplitud
          absoluta del ruido (siempre reemplaza por la mediana local).
    \item El kernel gaussiano con \(\sigma=0{,}8\) suaviza ruido de alta frecuencia
          en un rango universalmente aceptado por la literatura para imágenes
          fotográficas de resolución típica.
\end{itemize}

Los valores por defecto de la Tabla~\ref{tab:params} son constantes del código y no
se modificaron para ninguna imagen del dataset.

\begin{table}[H]
\centering
\caption{Parámetros del pipeline (constantes, no se ajustan por imagen).}
\label{tab:params}
\begin{tabular}{@{}lll@{}}
\toprule
Parámetro & Valor & Significado \\
\midrule
\texttt{gaussianSigma}   & \(0{,}8\)    & Desviación del kernel gaussiano \\
\texttt{medianSize}      & \([3,3]\)    & Tamaño del kernel de mediana \\
\texttt{claheClipLimit}  & \(0{,}005\)  & Límite relativo para recorte CLAHE \\
\texttt{claheTiles}      & \([8,8]\)    & Grilla de tiles CLAHE \\
\texttt{whiteBalance}    & \(\text{true}\) & Activa balance de blancos gray-world \\
\bottomrule
\end{tabular}
\end{table}

% ============================================================
\section{Implementación}
\label{sec:impl}

El código completo de la función se presenta a continuación. La función acepta
opcionalmente un \texttt{struct} de opciones, pero el demo del proyecto nunca lo
invoca, garantizando el carácter automático descrito en la
Sección~\ref{sec:auto}.

\begin{lstlisting}[caption={\texttt{src/preprocessing/pipeline\_preproc.m} — función principal.}]
function imgOut = pipeline_preproc(imgIn, opts)
% PIPELINE_PREPROC  Preprocesamiento automatico de imagenes de escritorio.

    % ---------- Defaults ----------
    if nargin < 2 || isempty(opts)
        opts = struct();
    end
    defaults = struct( ...
        'gaussianSigma',  0.8, ...
        'medianSize',     [3 3], ...
        'claheClipLimit', 0.005, ...
        'claheTiles',     [8 8], ...
        'whiteBalance',   true ...
    );
    opts = merge_opts(defaults, opts);

    % ---------- 1. Normalizacion a double [0, 1] ----------
    if isa(imgIn, 'uint8')
        img = im2double(imgIn);
    else
        img = imgIn;
    end

    % ---------- 2. Denoise: mediana en cada canal ----------
    img = cat(3, ...
        medfilt2(img(:,:,1), opts.medianSize), ...
        medfilt2(img(:,:,2), opts.medianSize), ...
        medfilt2(img(:,:,3), opts.medianSize));

    % ---------- 3. Denoise suave: gaussiano ----------
    img = imgaussfilt(img, opts.gaussianSigma);

    % ---------- 4. White balance gray-world ----------
    if opts.whiteBalance
        img = gray_world_wb(img);
    end

    % ---------- 5. CLAHE en canal L de Lab ----------
    lab = rgb2lab(img);
    L = lab(:,:,1) / 100;
    L = adapthisteq(L, ...
        'ClipLimit',   opts.claheClipLimit, ...
        'NumTiles',    opts.claheTiles, ...
        'Distribution','uniform');
    lab(:,:,1) = L * 100;
    img = lab2rgb(lab);

    img = max(0, min(1, img));
    imgOut = img;
end
\end{lstlisting}

\begin{lstlisting}[caption={Función auxiliar \texttt{gray\_world\_wb} — balance de blancos sin parámetros.}]
function out = gray_world_wb(img)
    r = img(:,:,1); g = img(:,:,2); b = img(:,:,3);
    meanR = mean(r(:)); meanG = mean(g(:)); meanB = mean(b(:));
    meanAll = (meanR + meanG + meanB) / 3;
    r = r * (meanAll / max(meanR, eps));
    g = g * (meanAll / max(meanG, eps));
    b = b * (meanAll / max(meanB, eps));
    out = cat(3, r, g, b);
    out = max(0, min(1, out));
end
\end{lstlisting}

% ============================================================
\section{Resultados experimentales}
\label{sec:resultados}

\subsection{Dataset y protocolo}

Se evaluó el pipeline sobre un conjunto de \(17\) imágenes de escritorio
(\texttt{dataset/reference/}) capturadas bajo condiciones de iluminación variadas:
luz natural diurna, lámparas de escritorio LED cálidas, y luz fluorescente de techo.
Las resoluciones oscilan entre \(853 \times 1280\) y \(1920 \times 1280\) píxeles.
El script \texttt{demo\_24abril.m} itera automáticamente sobre las imágenes, aplica el
pipeline y genera figuras comparativas.

\subsection{Métricas}

Para cada imagen se reportan PSNR y SSIM entre la imagen original (normalizada a
\([0,1]\)) y la salida del pipeline. La Tabla~\ref{tab:metricas} muestra los
resultados.

\begin{table}[H]
\centering
\caption{PSNR y SSIM entre entrada y salida del pipeline sobre el conjunto de 17 imágenes.}
\label{tab:metricas}
\begin{tabular}{@{}lcc@{}}
\toprule
Imagen & PSNR (dB) & SSIM \\
\midrule
08\_notebook\_pens\_pencils & 14.30 & 0.497 \\
09\_pen\_calculator         & 21.76 & 0.841 \\
p01\_desk\_pencil\_eraser   & 24.04 & 0.959 \\
p02\_pencils\_desk          & 24.27 & 0.962 \\
p03\_pencils\_close         & 25.35 & 0.958 \\
p04\_pen\_paper             & 23.67 & 0.934 \\
p05\_stationery             & 19.95 & 0.341 \\
p06\_stationery2            & 29.40 & 0.964 \\
p07\_stationery3            & 26.48 & 0.939 \\
p08\_stationery4            & 13.66 & 0.599 \\
p09\_markers\_desk          & 21.36 & 0.412 \\
p10\_markers\_colored       & 20.36 & 0.693 \\
p11\_markers2               & 21.27 & 0.377 \\
p12\_markers\_holder        & 22.88 & 0.353 \\
p13\_pens\_markers          & 24.15 & 0.891 \\
p14\_stationery\_flat       & 26.85 & 0.945 \\
p15\_pencil\_eraser2        & 23.63 & 0.940 \\
\midrule
\textbf{Promedio}           & \textbf{22.55} & \textbf{0.741} \\
\bottomrule
\end{tabular}
\end{table}

\subsection{Interpretación correcta de las métricas}

Debe subrayarse que, en este caso, PSNR y SSIM se calculan entre la imagen
\textbf{original} y la \textbf{salida} del pipeline; no existe un \emph{ground truth}
limpio con el cual compararlos. Por lo tanto:

\begin{itemize}
    \item Un PSNR alto significa ``el pipeline modificó poco la imagen''.
    \item Un PSNR bajo significa ``el pipeline introdujo muchos cambios'', lo cual
          es deseable precisamente en imágenes con problemas severos de iluminación
          o color.
\end{itemize}

En consecuencia, los valores bajos de PSNR observados en imágenes como
\texttt{p08\_stationery4} (13.66\,dB) o \texttt{p05\_stationery} (19.95\,dB)
\emph{no indican fallos} del pipeline: indican que esas imágenes eran aquellas en las
que el pipeline corrigió más intensidad y tinte. La validación cualitativa se realiza
visualmente mediante las figuras comparativas generadas.

\subsection{Inspección visual}

Las 17 figuras comparativas (original vs. preprocesada, lado a lado) se almacenan en
\texttt{results/figures/avance\_24abril/}. La Figura~\ref{fig:ejemplos} muestra dos
casos representativos: uno de alta similitud (poca corrección necesaria) y uno de
baja similitud (corrección agresiva de iluminación).

\begin{figure}[H]
\centering
\begin{subfigure}{\textwidth}
    \centering
    \includegraphics[width=0.95\linewidth]{../../results/figures/avance_24abril/p06_stationery2_comparacion.png}
    \caption{\texttt{p06\_stationery2} — PSNR 29.40, SSIM 0.964. Imagen bien iluminada; el pipeline aplica correcciones mínimas.}
\end{subfigure}
\vspace{0.5cm}
\begin{subfigure}{\textwidth}
    \centering
    \includegraphics[width=0.95\linewidth]{../../results/figures/avance_24abril/p08_stationery4_comparacion.png}
    \caption{\texttt{p08\_stationery4} — PSNR 13.66, SSIM 0.599. Imagen con fuerte dominante cromática; el pipeline redistribuye intensidades y corrige el balance.}
\end{subfigure}
\caption{Ejemplos representativos del efecto del pipeline.}
\label{fig:ejemplos}
\end{figure}

% ============================================================
\section{Discusión}

\subsection{Fortalezas}
\begin{itemize}
    \item Automaticidad real: ningún parámetro se ajusta por imagen; los operadores
          son adaptativos por construcción.
    \item Separación clara de responsabilidades: cada etapa ataca un problema
          específico (ruido impulsivo, ruido gaussiano, tinte cromático, iluminación no uniforme).
    \item Orden justificado: las dependencias entre etapas dictan una única
          secuencia razonable.
    \item Bajo costo computacional: en las pruebas realizadas, el pipeline procesa
          imágenes de \(\sim 1{,}5\) Mpíxeles en aproximadamente \(2\) segundos en un
          equipo estándar.
\end{itemize}

\subsection{Limitaciones y trabajo futuro}
\begin{itemize}
    \item La hipótesis del mundo gris puede fallar si la escena está dominada
          por un solo color. Una evolución razonable es sustituirla por el algoritmo
          \emph{Shades of Gray} o por \emph{max-RGB}, ambos igualmente automáticos.
    \item CLAHE sobre un único canal no corrige artefactos de color derivados de
          sub-exposición extrema. Una mejora consiste en aplicar tone-mapping antes
          del balance de blancos.
    \item La métrica PSNR/SSIM entre entrada y salida no mide calidad real.
          Una evaluación más rigurosa se obtendría midiendo el impacto del pipeline
          en la tasa de detección de la Parte~2 del proyecto.
\end{itemize}

% ============================================================
\section{Conclusión}

Se desarrolló un pipeline de preprocesamiento de imágenes de escritorio compuesto
por cinco etapas — normalización, filtrado de mediana, suavizado gaussiano, balance
de blancos \emph{gray-world} y normalización de intensidad mediante CLAHE sobre el
canal de luminancia — cuyo diseño cumple el requisito clave de automaticidad
exigido por el enunciado. Cada etapa se justifica en términos de un problema
concreto de la imagen (ruido impulsivo, ruido residual, dominante cromática,
iluminación no uniforme) y el orden de aplicación responde a las dependencias
causales entre ellas. La evaluación sobre 17 imágenes tomadas bajo condiciones
heterogéneas muestra que el pipeline preserva la estructura de las imágenes bien
capturadas (alta similitud) y corrige activamente aquellas con problemas de
iluminación o tinte (baja similitud), comportamiento coherente con lo esperado
para un algoritmo adaptativo.

% ============================================================
\appendix
\section{Reproducibilidad}

\begin{itemize}
    \item Sistema operativo: Ubuntu 22.04.
    \item MATLAB: R2026a (\(26.1.0.3203278\)).
    \item Toolboxes necesarios: \textit{Image Processing Toolbox}.
    \item Script de entrada: \texttt{demo\_24abril.m} (en la raíz del proyecto).
    \item Ejecución: \texttt{matlab -batch "demo\_24abril"} desde la carpeta del proyecto,
          o doble-clic y F5 desde el IDE.
    \item Las figuras comparativas se generan automáticamente en
          \texttt{results/figures/avance\_24abril/}.
\end{itemize}

\end{document}
